\section{Profile rigidity}\label{sec:profile}

We prove the two profile statements needed for propagation. At scales
tending to zero, the one-site closure and the comparison's first
variation force a linear rescaled overlap quantile. At scales bounded
away from zero, the optimizer converges to the scalar Parisi minimizer.
The latter gives the macroscopic starting profile; the former controls
the smaller scales reached by propagation.

Throughout this section, $T>1$ is fixed. All sequences have $n\to\infty$,
and the constants in the hypotheses are fixed before this limit is taken.
The paths to which a pressure or a planted expectation is applied are
nonnegative and nondecreasing. A bare endpoint may be used merely to
parameterize an admissible comparison on its working mass interval; no
pressure at an inadmissible bare endpoint is used below.

We first consider a small direct test. More precisely, the data satisfy
\eqref{eq:3}--\eqref{eq:5}, the path $r=r_n$ minimizes \eqref{eq:6} at a
specified positive strength, and $x=x_n\to0$ satisfies
\[
 x\gtrsim a,\qquad e(x)x^2\gtrsim nE^2,\qquad e(x)\gtrsim x^2.
\]
There are fixed constants $C_0,c_0,C_1>0$ such that
\[
 c_0u\le r(u)\le C_1u\qquad (C_0x\le u<1).
\]
The omitted mass, if present, is $u_0=o(x)$, and the initial radius obeys
\eqref{eq:5}. Total matrix and field clocks are bounded by fixed constants.
These are precisely the direct-test hypotheses of
Proposition~\ref{comp:direct-moments}. In particular, they imply $x=O(d)$.
Write $b=b_0-e(r)$, $h=h_0+e(r)r$ on the working interval and retain the
specified changed paths below it. Put
\[
 K_u=h_u+b_uS_u.
\]
This is a nonnegative nondecreasing scalar covariance path.

\begin{theorem}[Rigidity at a small test scale]\label{prof:direct-profile}
Under the preceding hypotheses, every sequence has a subsequence on
which, for some $\kappa$,
\begin{equation}
 \frac{r(xv)}{x}\longrightarrow\frac{v}{\kappa}
 \quad\text{locally uniformly for }v\in(0,\infty),
 \qquad 0<c\le\kappa\le C<\infty.
 \tag{28}\label{eq:28}
\end{equation}
The bounds on $\kappa$ depend only on $T$ and the total-clock bounds.
The same conclusion holds for a comparison with an omitted lower group,
provided its changed paths and the constrained competitors used in the
first variation are admissible.
\end{theorem}

\subsection{The closure estimates in both coordinates}

Physical mass $u$ parametrizes the planted splits, whereas overlap
radius $r$ parametrizes the inverse quantile $\alpha$ and its first
variation. We will take limits in both coordinates: a floor interval
can have negligible mass and nonnegligible radius length, so only the
radius estimates can detect a possible gap there. We first collect the
needed consequences of \eqref{eq:19} and \eqref{eq:22}.
For fixed $0<c<C<\infty$, let
\[
 I_n(c,C)=\{r:0\le r\le Cx,\ cx\le\alpha(r)\le Cx\}.
\]
At endpoints either fixed side convention may be used. We write $S(r)$,
$D(r)$, and $K(r)$ for the quantities at the corresponding split.

\begin{lemma}[Closure in mass and radius]\label{prof:coordinate-closure}
On each such band,
\[
 \frac1{x^3}\int_{I_n(c,C)}D(r)^2\,dr\longrightarrow0.
\]
The analogous statement holds with physical mass in place of radius.
Moreover, with
\[
 \rho_n(i)=\frac{S_i-\Gamma_n(K_i)}{x^3},
\]
there is a uniform bound $|\rho_n(i)|\le C_{c,C}$ on the band, and
\[
 \frac1x\int_{I_n(c,C)}|\rho_n(r)|\,dr\longrightarrow0.
\]
The corresponding normalized mass integral also tends to zero.
The uniform bound includes arbitrary moving split locations in the
specified bands; along any such sequence with $D_i=o(x)$,
$\rho_n(i)$ tends to zero.
\end{lemma}

\begin{proof}
The concentration assertions are \eqref{eq:19}; their radius version is
a direct estimate, rather than a change of variables in a mass estimate.
For clarity, the estimate that pays the inverse-density loss is
\[
 \frac{\operatorname{Var}(x^{-1}\int_I L\,dr)}{x^2}
 \lesssim\frac{n^{o(1)}}{ne(x)x^4k(Cx)}
 \lesssim\frac{n^{1-\delta+o(1)}E^2}{e(x)x^2}
 \lesssim n^{-\delta+o(1)}
\]
for radius subintervals of the band. Here one applies \eqref{eq:1} with
mass weight $\mathbf1_I(r(u))/(x\alpha'(r(u)))$, and uses the clock
inequality, $H_u\ge cxC_u$, and
$k(Cx)\gtrsim n^\delta u_*^2/x^2$. The ordered side-average argument
in Proposition~\ref{comp:direct-moments} converts this estimate and the
integrated trace estimate to the asserted radius mean-square bound.
Thus the conclusion remains valid even when $\alpha'=k$ throughout
most of the radius interval.

The raw bounds \eqref{eq:12} give $D_i=O(x)$ uniformly here.
Proposition~\ref{prop:cavity-closure} supplies a fixed $\vartheta>0$
and $\eta_n\to0$ such that
\[
 |\rho_n(i)|\le C\{\eta_n+(D_i/x)^\vartheta\}
\]
at every split in the band. This gives the uniform bound on $\rho_n$,
including moving splits. For either normalized measure $du/x$ or
$dr/x$, the band has bounded measure and $D/x$ tends to zero in
$L^2$ while remaining uniformly bounded. H\"older's inequality
therefore makes the integral of $(D/x)^\vartheta$ tend to zero.
The displayed modulus proves both asserted $L^1$ limits, including
the radius limit on floor intervals.
\end{proof}

\subsection{Scalar derivatives and rescaled compactness}

Fix a full scalar covariance path $K=(K_u)$. Write $\Gamma(k)$ for
its scalar moment curve at covariance value $k$; differentiation in $k$
holds that path, and hence its mass profile in scalar time, fixed.
The subscript $n$ records the full path of the $n$th system.
In scalar time $q=k/T$, put
\[
 \alpha_K(q)=|\{u\in(0,1):K_u/T\le q\}|.
\]
A jump of $K$ is filled with its constant mass, and an initial empty
clock is filled with mass zero. The scalar continuation $V$ and its
optimal diffusion satisfy
\[
 V_q=-\frac T2\bigl(V_{zz}+\alpha_K(q)V_z^2\bigr),\qquad
 dZ_q=T\alpha_K(q)V_z(q,Z_q)\,dq+\sqrt T\,dB_q,
 \qquad Z_0=0.
\]
The terminal function is $\log\cosh z$. Adding a terminal mass-one
clock changes $V$ only by a constant in $z$ and does not change its
positive-order spatial derivatives. We use
\[
 \widehat\Gamma(q)=\Gamma(Tq)
       =\mathbb E V_z(q,Z_q)^2.
\]

The moment identities below are the SK specialization of
\cite[Proposition~3]{AuffingerChenProperties2015}; see also
\cite[Lemma~8.4]{JagannathTobasco2017}. We give their diffusion
derivation together with the uniform fourth-derivative bound needed
for the limiting slope.

\begin{lemma}[Scalar derivatives and strict curvature]
\label{prof:scalar-derivatives}
For scalar clocks with a fixed bound on total variance, positive-order
spatial derivatives of every fixed order are uniformly bounded.
They are absolutely continuous in scalar time, with uniformly bounded
almost-everywhere time derivatives of each fixed spatial order.
The function $V$ is even in $z$, $V_{zz}>0$, and
\[
 \widehat\Gamma'(q)=T\mathbb E V_{zz}(q,Z_q)^2,
 \qquad
 \widehat\Gamma''(q)=T^2\mathbb E\bigl[
 V_{zzz}(q,Z_q)^2-2\alpha_K(q)V_{zz}(q,Z_q)^3\bigr]
\]
for almost every $q$. There are constants $c,C>0$, depending only on
$T$ and the total-clock bound, such that
\[
 c\le V_{zz}(0,0)\le C,
 \qquad c\le -V_{zzzz}(0,0)\le C.
\]
\end{lemma}

\begin{proof}
Corollary~\ref{cor:found-source-bounds}, applied to one Ising spin,
gives the spatial derivative bounds independently of the partition.
The finite-step formulas pass to general clocks by
Proposition~\ref{prop:found-limits}.
The PDE expresses the time derivatives in terms of bounded spatial
derivatives and $0\le\alpha_K\le1$. Convexity and evenness are preserved
by every scalar Gaussian continuation.

It\^o's formula gives
\[
 dV_z(q,Z_q)=\sqrt T\,V_{zz}(q,Z_q)\,dB_q,
\]
and
\[
 dV_{zz}(q,Z_q)=-T\alpha_K(q)V_{zz}(q,Z_q)^2\,dq
                 +\sqrt T\,V_{zzz}(q,Z_q)\,dB_q.
\]
Taking expectations of the squares proves the two displayed derivative
identities. Absolute continuity is sufficient for these applications
of It\^o's formula; alternatively one first works with step clocks and
passes to the limit using the stated derivative bounds.

For the strict lower bounds, let $Q$ be terminal scalar time.
Differentiating the equation once more shows that $w=V_{zz}$ solves
a linear backward equation with drift $T\alpha_KV_z$, potential
$T\alpha_Kw$, and terminal value $\operatorname{sech}^2z$.
Since the drift is uniformly bounded and the potential is nonnegative,
its Feynman--Kac representation bounds $w(0,0)$ below by a positive
constant uniformly for bounded $Q$. One can see the uniformity directly
by restricting the terminal diffusion to a fixed bounded interval,
which has uniformly positive probability.

Set $v=V_{zzz}$. It satisfies the linear equation
\[
 v_q+\frac T2v_{zz}+T\alpha_KV_zv_z
          +3T\alpha_KV_{zz}v=0.
\]
On $z>0$, the drift and potential are nonnegative, $v(q,0)=0$, and
the terminal value is $-2\tanh z\,\operatorname{sech}^2z<0$.
Use the killed-diffusion representation on the positive half-line.
For $Q$ in a fixed compact subinterval of $(0,\infty)$ and small
initial $z>0$, the event that $z+\sqrt T B$ stays positive and ends
in a fixed interval $[a,b]\subset(0,\infty)$ has probability at least
$c_1z$, by the reflected Gaussian kernel. On this event the diffusion
with nonnegative bounded drift also survives, and its terminal value
lies in $[a,b+TQ]$. Its negative terminal datum is therefore bounded
away from zero in absolute value, while the Feynman--Kac weight is at
least one. Hence $-v(0,z)\ge c_2z$, and division by $z$ gives
$-V_{zzzz}(0,0)\ge c_2$. If $Q$ tends to zero, the bounded time
derivative and $V_{zzzz}(Q,0)=-2$ give the same conclusion. Upper
bounds follow from the spatial derivative bounds.
\end{proof}

We next compare the overlap radius with scalar time $q=K/T$.
At the test scale their distributions have the same limit. The scalar
second derivative then has size $x$ on clocks of size $x$; two
integrations produce the cubic defect below once its linear term has
been calibrated at a positive outer split.

\begin{lemma}[Scalar compactness at the test scale]
\label{prof:scalar-compactness}
Under the hypotheses of Theorem~\ref{prof:direct-profile}, pass to a
subsequence. There is a locally finite measure $\nu$ on $[0,\infty)$,
with cumulative function $U(t)=\nu([0,t])$, such that the rescaled
distributions of $r$ and of $K/T$ both converge locally weakly to $\nu$.
For sufficiently large $t$, $U(t)$ is bounded above and below by fixed
positive multiples of $t$. Further subselection gives
\begin{equation}
 \begin{gathered}
 g_n(t)=\frac{\Gamma_n(Txt)-xt}{x^3}
       \longrightarrow g(t)\quad\text{in }C^1_{\mathrm{loc}}([0,\infty)),
 \\
 g(0)=0,\qquad g''(t)=At-BU(t)\quad\text{for almost every }t,
 \qquad A,B>0.
 \end{gathered}
 \tag{29}\label{eq:29}
\end{equation}
The function $g$ vanishes on $\operatorname{supp}\nu$. The constants
$A,B$ have positive lower bounds and finite upper bounds depending
only on $T$ and the total-clock bound.
\end{lemma}

\begin{proof}
\emph{The limiting mass distribution.}
Give a Borel radius set $J$ rescaled mass
$x^{-1}|\{u>u_0:r(u)/x\in J\}|$. Its cumulative function differs from
$\alpha(xt)/x$ by at most $u_0/x=o(1)$. Outer comparability bounds its
mass on each compact set and bounds its cumulative function above
and below by multiples of $t$ for sufficiently large $t$. Helly
selection gives a locally finite weak limit $\nu$.

We next compare the scalar clock. The total floor mass is bounded by
$\int k(r)\,dr=o(x)$. On a positive compact mass band, outer
comparability and monotonicity give $r=O(x)$, while cap widths and the
initial wall are $o(x)$. Since $x=O(d)$ and $\lambda\le d^2$, the
taper argument $\lambda r/d^3$ is bounded on these bands, including
their cap endpoints. Thus $1-p$ has a fixed positive lower bound.
The proof of Lemma~\ref{comp:geometry} then gives \eqref{eq:10}
with a fixed constant in place of $n^{o(1)}$ on these bands.
Together with \eqref{eq:9} and the mass estimate in \eqref{eq:19},
this proves $S-r=o(x)$ in normalized mass measure off a negligible
set. Since $b_0-T=o(x^2)$, $h_0=o(x^3)$, and $S=O(x)$ there,
\[
 K/T-r=(1-e/T)(S-r)+((b_0-T)S+h_0)/T=o(x)
\]
in the same measure. Mass below $\varepsilon x$ contributes at most
$\varepsilon$ to the rescaled measure. Monotonicity and an outer
regular point at a sufficiently large fixed multiple of $x$ prevent
mass above a sufficiently large multiple of $x$ from contributing
to a fixed compact clock interval. Letting $\varepsilon\downarrow0$
proves the same local weak limit for $K/(Tx)$, including possible
atoms at zero. In particular $\alpha_K(xt)/x$ is locally bounded
and converges to $U(t)$ at its continuity points.

\emph{The second derivative of the scalar defect.}
Write $a_{2,n}=V_{n,zz}(0,0)$ and $a_{4,n}=V_{n,zzzz}(0,0)$.
Select convergent subsequences of these two numbers. On compact
$t$-intervals the diffusion has
\[
 \frac{\mathbb E Z_{xt}^2}{x}\longrightarrow Tt,
 \qquad \mathbb E|Z_{xt}|^{2j}=O_j(x^j).
\]
These statements follow from bounded drift, or directly from its
bound $T\alpha_K|V_z|\le Cx|Z|$ on the relevant clocks. Oddness
and the derivative bounds give
\[
 \frac{\mathbb E V_{n,zzz}(xt,Z_{xt})^2}{x}
       \longrightarrow T a_4^2t,
 \qquad
 \mathbb E V_{n,zz}(xt,Z_{xt})^3\longrightarrow a_2^3.
\]
For example,
$V_{n,zzz}(q,z)=a_{4,n}z+O(q|z|+|z|^3)$ uniformly on bounded
$z$-intervals, and the moment bounds control the complement.
By Lemma~\ref{prof:scalar-derivatives},
\[
 \begin{aligned}
 g_n''(t)&=\frac{\widehat\Gamma_n''(xt)}x\\
 &=T^2\left[
 \frac{\mathbb E V_{n,zzz}(xt,Z_{xt})^2}{x}
 -2\frac{\alpha_K(xt)}x\mathbb E V_{n,zz}(xt,Z_{xt})^3\right].
 \end{aligned}
\]
These second derivatives are locally bounded and converge almost
everywhere, and in local $L^1$, to
\[
 At-BU(t),\qquad A=T^3a_4^2,\qquad B=2T^2a_2^3.
\]
Both coefficients have the asserted strict bounds.

\emph{Calibration of the linear term.}
It remains to bound $g_n'(0)$; a bound on $g_n''$ alone would leave
this coefficient undetermined. Choose a nonfloor split in an
outer band with $u,r\asymp x$ and $D=o(x)$, using the mass version
of \eqref{eq:19}. Its scalar clock $K/(Tx)$ is bounded above and
below by positive constants. If $e(x)/x^2$ is bounded, then
$e(0)=O(x^2)$ on these compact radius bands and
$(S-K/T)/x^3=O(1)$. If $e(x)/x^2\to\infty$, then $x/d\to0$ and
$e=\lambda$, $p=0$ exactly on every compact radius band. On contact
$S=r$; on a cap,
\[
 |S-r|\le \ell+r_{\min}\mathbf1_{\mathrm{initial}},\qquad
 \frac{\lambda(\ell+r_{\min})}{x^3}
 \lesssim n^{-\delta}\log n\,
          \frac{\lambda}{d^2}\left(\frac ax\right)^3=o(1).
\]
Here we used $L\le a^3/d^2$, the cap-width bound, and
$r_{\min}\le n^{-\delta}L$. Thus the same scalar value is bounded
also in this regime. Closure \eqref{eq:22} bounds $g_n$ at this
positive clock. As $g_n(0)=0$ and $g_n''$ is locally bounded, it
bounds $g_n'(0)$. Subsequence selection and integration of $g_n''$
prove \eqref{eq:29}.

\emph{Zeros on the support.}
The same argument at almost every selected mass point gives $g=0$
on the support of $\nu$. More explicitly, in the bounded-strength
regime $e(S-r)=o(x^3)$ on the full-measure selections above; in the
constant-profile regime the displayed cap estimate gives this
conclusion. The bare-clock error is $o(x^3)$, and
Lemma~\ref{prof:coordinate-closure} supplies little-oh closure at
such points. Therefore $g_n(K/(Tx))\to0$ on full normalized-mass
selections. Every positive support point can be approached by these
selections from neighborhoods carrying positive limiting mass,
including an approximating atom. Local uniform convergence of $g_n$
then gives $g=0$ there. At zero use $g(0)=0$.
\end{proof}

\Needspace{3\baselineskip}
\begin{lemma}[Uniform scalar linearization]
\label{prof:low-clock-linearization}
For each fixed $C<\infty$, under the small direct-test hypotheses,
\[
 \sup_{0\le K\le Cx}|\Gamma_n(K)-K/T|=O(x^3),
 \qquad \Gamma'_{n,K}(0)=T^{-1}+O(x^2).
\]
Only the portion of this clock interval present in the scalar
continuation is needed; for all sufficiently large $n$ it contains
every displayed fixed multiple of $x$. The assertion is uniform
along arbitrary admissible sequences with the stated fixed bounds.
It allows an omitted mass $u_0=o(x)$ and requires no version of
\eqref{eq:4} on that omitted group.
\end{lemma}

\begin{proof}
Choose the outer calibration band in the preceding proof beyond the
fixed target interval. At its selected split, $K_c\asymp x$ and
$K_c>Cx$. Monotonicity gives $\alpha_K(q)=O(x)$ for every smaller
clock; the omitted mass contributes only $o(x)$. The second-derivative
identity then gives $\widehat\Gamma''(q)=O(x)$ for $q=O(x)$.
The positive-clock calibration gives
$\Gamma(K_c)-K_c/T=O(x^3)$, and $\Gamma(0)=0$. Taylor integration
therefore yields the derivative estimate and the uniform bound down
to zero, including clock gaps. Equivalently, they follow from the
uniform local bounds on $g_n$ and $g_n'$ in the preceding proof.
If uniformity along the allowed sequences failed, selecting the
violating sequence and repeating that proof would be a contradiction.
No lower-group stationarity or lower-group closure is used here.
\end{proof}

\subsection{Virtual gaps and limiting potentials}

The scalar defect $g$ now vanishes on the support of $\nu$. To show
that the support has no gaps, we use the first variation at radii
that may carry no limiting mass. This is where the radius form of
the closure estimate is needed.

\begin{lemma}[Stationarity on virtual radius intervals]
\label{prof:virtual-potential}
Let $l=\inf\operatorname{supp}\nu$. On compact radius intervals
$t>l$, subselection gives a monotone limit
$s(t)=\lim S(xt)/x$ almost everywhere. The normalized stationarity
quantity $G(xt)/x^3$ is uniformly bounded on each such interval and
converges strongly in $L^1$ to a function $F$. There are two cases.
If $e(x)/x^2\to\infty$, then
\[
 s(t)=t,\qquad F(t)=-2Tg(t).
\]
Otherwise, after subselection, $e(xt)/x^2\to h_1(t)>0$ and the
logarithmic slopes converge locally; then
\[
 g(s(t))=\frac{h_1(t)}T(s(t)-t),
\]
and $F(t)$ has the sign of $t-s(t)$, with equality precisely when
$t=s(t)$.

If $(a_1,a_2)$ is a gap in $\operatorname{supp}\nu$ with endpoints
in that support and positive cumulative mass on the gap, there is
a nonnegative Lipschitz function $\mathcal H$ on $[a_1,a_2]$ with
zero endpoint values such that
\[
 \mathcal H'(t)=j(t)F(t)\quad\text{almost everywhere},
 \qquad 0<c\le j(t)\le C.
\]
In the bounded-strength case, $s(t)\in[a_1,a_2]$ almost everywhere
on the gap.
\end{lemma}

\begin{proof}
On each compact interval above $l$, masses $\alpha(xt)$ lie between
fixed positive multiples of $x$. The raw estimates give $S=O(x)$,
$B=O(x^2)$, and $K=O(x)$ there. Helly selection supplies $s$.
Introduce the bare-clock remainder
\[
 \zeta_n(r)=\frac{(b_0(\alpha(r))-T)S(r)+h_0(\alpha(r))}{x^3}.
\]
By \eqref{eq:4}, it is uniformly $o(1)$ on these compact mass bands.
The exact closure identity is
\[
 \frac{e(r)(S(r)-r)}{Tx^3}
   =g_n\!\left(\frac{K(r)}{Tx}\right)+\rho_n(r)+\frac{\zeta_n(r)}T.
\]
In particular the closure error is additive in this identity.
Lemma~\ref{prof:coordinate-closure} gives both its uniform bound
and its strong radius $L^1$ convergence to zero.

If $e(x)/x^2\to\infty$, then
$\lambda r/d^3\le Cx/d\to0$ on a compact radius band. Thus
$e=\lambda$ and $p=0$ exactly there. The last identity and the
boundedness of $g_n$ imply
\[
 \sup\frac{|S-r|}{x}\lesssim\frac{x^2}{e}\longrightarrow0,
 \qquad \frac{K(r)}{Tx}-\frac rx\longrightarrow0
\]
uniformly on the band. Therefore $G=2e(r-S)$ is uniformly of order
$x^3$ and $G(xt)/x^3\to-2Tg(t)$ strongly in $L^1$.

In the other case, the fixed smooth strength profile has a locally
convergent subsequence with positive limit $h_1$; the same is true
of its logarithmic slope. Now $e=O(x^2)$, so $K/(Tx)-S/x\to0$.
Pass to the limit in the exact closure identity to obtain the
displayed equation for $g(s)$. Since $B=S^2+D^2$ and the radius
version of \eqref{eq:19} makes $D^2/x^2$ tend to zero in $L^1$,
the limit of $G/x^3$ is
\[
 F(t)=h_1(t)\left((2-p(t))t-2(1-p(t))s(t)
                            -p(t)\frac{s(t)^2}{t}\right).
\]
The bracket factors as
\[
 (t-s(t))\left(2-p(t)+p(t)\frac{s(t)}t\right),
\]
and its second factor is positive. Raw bounds give a uniform bound
on $G/x^3$ on compact positive radius intervals. If a band extends
to zero, $p=0$ on a fixed scaled neighborhood of zero because
$x=O(d)$ and $\lambda\le d^2$; this removes the apparent $1/t$
singularity.

We construct the potential, including its endpoint values. On any
band with masses comparable to $x$, normalize the first-variation
factor $J$ at one reference radius. Its distortion is bounded above
and below: indeed
\[
 \int wM'(\alpha)\,dr
 \le\int \frac{M'(\alpha)}{M(\alpha)}\alpha'\,dr
 \le2\log\frac{\alpha_{\mathrm{right}}+u_*}
                    {\alpha_{\mathrm{left}}+u_*}.
\]
Thus $J/J_{\mathrm{ref}}$ has a weak-star convergent subsequence in
$L^\infty$ with limit $j$ bounded above and below by positive
constants. Normalize the potential by $x^4J_{\mathrm{ref}}$.
Its derivative in scaled radius is $(J/J_{\mathrm{ref}})G/x^3$,
which is uniformly bounded. Strong $L^1$ convergence of $G/x^3$
identifies the derivative of every uniform potential limit as $jF$.

A support point is approached by nonfloor mass points, since a
neighborhood with positive limiting rescaled mass cannot consist
of the $o(x)$ total floor mass. At contact the potential is zero.
At a point in a negative cap its right zero lies at radius distance
$o(x)$; the derivative bound makes the normalized potential there
$o(1)$. Outer comparability bounds the masses above by $Cx$ after
a fixed enlargement of the band. If the point is the first support
point and begins a gap, it must be an atom. Choose the reference
points inside its approximating atom after a fixed positive amount
of rescaled mass has accumulated. The mass is then bounded below
by $cx$ all the way from that reference to the interior of the gap;
no estimate on smaller masses is needed. The normalized potentials
are controlled on intervals starting at these moving reference
points. Their uniform Lipschitz bound extends the limit to the atom
with endpoint value zero. This also covers an atom at radius zero,
without requiring a bound at the exact prelimit radius zero.

These reference points both bound the additive constant of the
normalized potential and pin its limiting endpoint values to zero.
Negative cap components have radius width $o(x)$, so their negative
excursions disappear under the same Lipschitz bound. The limit is
therefore nonnegative. This uses small floor mass only to select a
reference point inside positive support mass, never to discard a
floor interval in radius.

Finally, in the bounded-strength case, monotonicity bounds the
virtual values of $S/x$ between active mass selections approaching
the two support endpoints. At those selections $S-r=o(x)$.
Consequently $a_1\le s(t)\le a_2$ on the gap, as claimed.
\end{proof}

\begin{proof}[Proof of Theorem~\ref{prof:direct-profile}]
We show that the limiting support is all of $[0,\infty)$. On a
bounded gap, the variational potential forces both endpoint slopes
of $g$ to be nonpositive, while its scalar equation forces their
sum to be positive. Consider such a nontrivial gap
$(a_1,a_2)$ with positive cumulative mass. We claim
$g'(a_1)\le0$ and $g'(a_2)\le0$.
If $g'(a_1)>0$, then $g>0$ immediately to the right of $a_1$.
In the diverging-strength case this makes $F<0$ there, contrary to
$\mathcal H\ge0$ and $\mathcal H(a_1)=0$. In the bounded-strength
case the closure equation excludes $s=a_1$ for $t>a_1$. If $s$ is
close to $a_1$, the inequality $g(s)>0$ forces $s>t$; if $s$ is
farther from $a_1$, it is again greater than $t$ for all sufficiently
nearby $t$. Thus $F<0$ almost everywhere immediately to the right,
giving the same contradiction. If $g'(a_2)>0$, the analogous
argument on the left gives $F>0$ there, which contradicts
$\mathcal H\ge0$ and $\mathcal H(a_2)=0$.

On the gap, $U$ is constant, $g(a_1)=g(a_2)=0$, and
$g''(t)=At-\text{constant}$. Direct integration gives
\[
 g'(a_1)+g'(a_2)=\frac A6(a_2-a_1)^2>0,
\]
a contradiction. The support is unbounded because $U(t)\asymp t$
for large $t$. Hence, with $l=\inf\operatorname{supp}\nu$, it has
no gaps on $(l,\infty)$. There $g=0$, so \eqref{eq:29} implies
$U(t)=(A/B)t$, first almost everywhere and then everywhere by
right continuity. If $l>0$, then $g(l)=g'(l)=0$ and $g''(t)=At$
for $0<t<l$. These imply $g(0)=Al^3/3>0$, a contradiction.
Thus $l=0$ and $U(t)=\kappa t$, where $\kappa=A/B$.

This cumulative function is continuous and strictly increasing.
Local weak convergence of the rescaled measures therefore implies
convergence of their inverse quantiles at every positive mass,
and monotonicity makes this convergence locally uniform. The
strict bounds on $A,B$ give the bounds on $\kappa$. This proves
\eqref{eq:28}.
\end{proof}

\subsection{The macroscopic variational limit}

The preceding theorem handles scales tending to zero. For the
macroscopic starting point of propagation we also need to identify
the optimizer when the test scale stays bounded away from zero.
We first pass an optimized interpolation to its scalar limit, then
identify that limit's minimizer.

Let $f_{\rm spin}(K)=f_n(0,K)$. Independence across sites makes this
quantity independent of $n$. Its normalization includes the field
self-correction $-K(1-)/2$. A quantile in this subsection means a
nondecreasing measurable function from $(0,1)$ to $[0,1]$; it need
not be onto or strictly increasing.

\begin{lemma}[Macroscopic optimized interpolation]
\label{prof:macro-variational}
Let $b,h$ be bounded nonnegative nondecreasing paths, with
$b(u)\ge b_*>0$. Then
\[
 \lim_{n\to\infty} f_n(b,h)
 =\min_z\left\{f_{\rm spin}(h+bz)
                         +\frac14\int_0^1 b(u)z(u)^2\,du\right\}.
\]
The conclusion also holds for uniformly bounded paths $b_n,h_n$
converging to $b,h$ in $L^1(0,1)$, provided every actual path is
admissible. No continuity of the limiting paths is required.
\end{lemma}

\begin{proof}
Write $\eta_n=n^{-1/8}$ and $Q_n=n^{1/8}$, and first minimize over
onto quantiles whose inverse $\alpha:[0,1]\to[0,1]$ satisfies
\[
 \eta_n\le\alpha'\le Q_n,\qquad \alpha(0)=0,\quad\alpha(1)=1.
\]
This is a nonempty compact class of inverse paths at fixed $n$.
For $0\le t\le1$, put
\[
 \Phi_n(t)=\min_z\left\{
 f_n((1-t)b,h+tbz)+\frac t4\int_0^1 bz^2\,du\right\}.
\]
All the interpolated paths are admissible. The first identity in
\eqref{eq:1} and the envelope argument of Lemma~\ref{comp:envelope}
give local Lipschitz continuity and, at almost every $t$, the
derivative formula at a minimizing path,
\[
 0\le\Phi_n'(t)=\frac14\int_0^1
                       b\bigl[D^2+(S-z)^2\bigr]\,du.
\]
An upper bound by the right side would also suffice. All these
derivatives are bounded uniformly in $n,t$.

For $t>0$, first variation in the quantile has bracket
$t b(u)(r-S)/2$, where $r=z(u)$. The inverse-density constraints
are now constant. The proof of Lemma~\ref{comp:geometry} gives a
potential with derivative a positive constant times $b(u)(r-S)$:
the inverse density equals $Q_n$ on its negative components and
$\eta_n$ on its positive components. On its zero set $S=r$
almost everywhere. A negative component has radius width at most
$Q_n^{-1}$, since its mass is at most one. Its endpoint signs and
monotonicity of $S$ place $S(r)$ between its radius endpoints. This
also holds for an initial or final negative component, using
$0\le S\le1$ at the exterior endpoint. Consequently
\[
 \int_0^1(S-z)^2\,du\le Q_n^{-2}+\eta_n.
\]
Indeed the total mass on the floor is at most $\eta_n$, and off the
floor the pointwise error is at most $Q_n^{-1}$.

We give a concentration argument uniform over these minimizers for
$t$ bounded away from zero. In radius coordinates the field clock
satisfies $\partial_r(h+tbz)\ge tb_*$ in the sense of measures;
the matrix clock is also nondecreasing. The second identity in
\eqref{eq:1}, with nonnegative jump contributions discarded, yields
\[
 \int_0^1\mathbb E\operatorname{tr}H_r^2\,dr
        \le\frac1{ntb_*},\qquad
 \int_0^1\mathbb E\operatorname{tr}H_u^2\,du
        \le\frac{Q_n}{ntb_*}.
\]
Fix $\varepsilon>0$. Since $H_u\ge uC_u$, the trace and
Brascamp--Lieb bounds in \eqref{eq:1} imply
\[
 \int_\varepsilon^1\mathbb E\operatorname{tr}C_u^2\,du
       \le\frac{Q_n}{ntb_*\varepsilon^2},\qquad
 \operatorname{Var}\left(\int_I L_u\,du\right)
       \le\frac{C Q_n}{ntb_*\varepsilon}
\]
for every interval $I\subset[\varepsilon,1]$. For the second bound
use $\mathbb E W_u=\mathbb E\operatorname{tr}(H_uC_u)
\le u^{-1}\mathbb E\operatorname{tr}H_u^2$.

Here are the details converting interval concentration to pointwise
concentration in integrated mass. Partition $[\varepsilon,1]$
into at most $1+\lceil\Delta^{-1}\rceil$ intervals on which the
nondecreasing deterministic function $S$ varies by at most
$\Delta$, allowing empty intervals. Remove a strip of mass
$\rho$ at each end of each interval, and discard any interval
shorter than $2\rho$. The discarded mass is at most
$C\rho(1+\Delta^{-1})$. At a retained mass $u$, both
$I_-=[u-\rho,u]$ and $I_+=[u,u+\rho]$ lie in the same interval.
The martingale property of $X$ gives, for $v\le u\le w$,
\[
 \mathbb E_u L_v
   =\frac{|X_u|^2}{2}-\frac{|X_u-X_v|^2}{2}
   \le\frac{|X_u|^2}{2}
   \le\frac{\mathbb E_u|X_w|^2}{2}
   =\mathbb E_u L_w.
\]
After averaging over $I_-$ and $I_+$ and subtracting $S_u/2$,
conditional Jensen's inequality and the interval variance bound give
\[
 \bigl\||X_u|^2-S_u\bigr\|_2
 \le C\Delta+\frac C\rho
                   \left(\frac{Q_n}{ntb_*\varepsilon}\right)^{1/2}.
\]
On discarded masses use $0\le|X_u|^2,S_u\le1$. Thus, first
letting $n\to\infty$ for fixed $t\ge\varepsilon,\Delta,\rho$,
then $\rho\downarrow0$ and $\Delta\downarrow0$, proves
\[
 \sup_{t\in[\varepsilon,1]}
 \int_\varepsilon^1\operatorname{Var}(|X_u|^2)\,du\longrightarrow0.
\]
This order of limits is legitimate because every bound above is
uniform over the minimizing paths and over $t\ge\varepsilon$.
The identity
\[
 D_u^2=\operatorname{Var}(|X_u|^2)
       +\mathbb E\operatorname{tr}C_u^2
       +2\mathbb E X_u^{\mathsf T}C_uX_u
\]
and $\mathbb E X_u^{\mathsf T}C_uX_u
\le(\mathbb E\operatorname{tr}C_u^2)^{1/2}$ now show that the
same integrated convergence holds for $D^2$. Masses $u<\varepsilon$
and interpolation times $t<\varepsilon$ cost $O(\varepsilon)$.
Sending $\varepsilon\downarrow0$ in the integrated envelope bound
proves $\Phi_n(1)-\Phi_n(0)\to0$.

At $t=0$ the value is $f_n(b,h)$, independently of $z$. At $t=1$
it is exactly the scalar objective in the statement, restricted to
the inverse-density class. Every probability distribution on $[0,1]$
can be approximated weakly by smooth strictly positive densities
bounded above and below; their quantiles converge in $L^1$.
Each such approximating density belongs to the constrained class
for all sufficiently large $n$. Conversely, monotone quantiles form
a compact set in $L^1$. The variation formula gives
\[
 |f_n(b,h)-f_n(\widetilde b,\widetilde h)|
 \le\frac14\|b-\widetilde b\|_1
      +\frac12\|h-\widetilde h\|_1
\]
between any two admissible pairs, uniformly in $n$; the same bound
holds for $f_{\rm spin}$. These continuity bounds pass both the
constrained minima and the varying endpoint paths to their limits.
They also prove attainment of the scalar minimum.
\end{proof}

\subsection{The scalar minimizer and its initial slope}

For a quantile $q$ let $\alpha$ be its distribution function, and
let $V_\alpha$ solve the scalar PDE with total time one and terminal
value $\log\cosh z$. Integration of a terminal mass-one clock and
the layer-cake identity give, respectively,
\[
 f_{\rm spin}(Tq)=V_\alpha(0,0)-\frac T2,
 \qquad
 \int_0^1q(u)^2\,du=1-2\int_0^1s\alpha(s)\,ds.
\]
Consequently
\[
 f_{\rm spin}(Tq)+\frac T4\int_0^1q(u)^2\,du
 =V_\alpha(0,0)-\frac T2\int_0^1s\alpha(s)\,ds-\frac T4.
\]
Thus our scalar objective has the same minimizing distributions as
the usual zero-field SK Parisi functional at $\beta=\sqrt T$.

\begin{lemma}[Uniqueness and initial density]
\label{prof:scalar-minimizer}
The scalar variational problem has a unique minimizing distribution,
denoted by $\alpha_T$, and a unique quantile up to null sets,
denoted by $q_T$. There is a constant $\kappa_T>0$ such that
\[
 \alpha_T(s)=\kappa_Ts+o(s)\quad(s\downarrow0).
\]
Its quantile is continuous at every mass in $(0,1)$ and satisfies
$c_Tu\le q_T(u)\le C_Tu$ there.
\end{lemma}

\begin{proof}
Existence follows from compactness and the continuity just proved.
We include the uniqueness argument to specify the needed form, using
the stochastic-control and strict-convexity method of
\cite[Theorems~2--3]{AuffingerChenUniqueness2015}.
The scalar control representation is
\[
 V_\alpha(0,0)=\sup_{|m|\le1}
 \mathbb E\left[
  \log\cosh\left(\sqrt T B_1+T\int_0^1\alpha(s)m_s\,ds\right)
    -\frac T2\int_0^1\alpha(s)m_s^2\,ds\right],
\]
where controls are progressively measurable. To verify it, apply
It\^o's formula to $V_\alpha(t,Z_t)$ for
$dZ_t=T\alpha(t)m_t\,dt+\sqrt T\,dB_t$. The resulting difference
between the initial PDE value and the displayed payoff is
$\frac T2\mathbb E\int\alpha(t)(m_t-V_{\alpha,z}(t,Z_t))^2\,dt$.
It vanishes for the bounded feedback control
$m_t=V_{\alpha,z}(t,Z_t)$.

Suppose $\alpha_1$ and $\alpha_2$ differ on a set of positive
Lebesgue measure. Use the optimal control $m$ for their midpoint
$\bar\alpha$ in all three functionals. For this control,
$m_0=0$ by symmetry and
\[
 dm_t=\sqrt T\,V_{\bar\alpha,zz}(t,Z_t)\,dB_t,
 \qquad V_{\bar\alpha,zz}(t,Z_t)>0.
\]
With $a=\alpha_1-\alpha_2$, stochastic Fubini gives
\[
 \int_0^1a(t)m_t\,dt
 =\sqrt T\int_0^1
       \left(\int_s^1a(t)\,dt\right)
                    V_{\bar\alpha,zz}(s,Z_s)\,dB_s.
\]
Its second moment is strictly positive by It\^o isometry: a
vanishing continuous primitive would imply $a=0$ almost everywhere.
Consequently the terminal arguments for $\alpha_1$ and $\alpha_2$
differ with positive probability. Strict convexity of
$\log\cosh$ makes the midpoint payoff strictly smaller than the
average of the two payoffs for this control, which are bounded by
the corresponding optimal values. The integral penalty in the
Parisi functional is linear in $\alpha$. Strict convexity of the full
functional therefore proves uniqueness. Distribution functions equal almost
everywhere describe the same probability measure on $[0,1]$.

Lopatto's support theorem \cite[Theorem~1.1]{Lopatto2026} states
that, at zero field and $T=\beta^2>1$, the minimizing measure has
support $[0,q_\beta]$ with $q_\beta>0$, no atom at zero, and a smooth
density near zero. Its scalar martingale satisfies
$\widehat\Gamma(s)=s$ on that support by
\cite[Lemma~2.2, equation~(2.3)]{Lopatto2026}.
We now compute the initial density and prove its positivity.
Since $\widehat\Gamma''(s)=0$ for $0<s<q_\beta$,
Lemma~\ref{prof:scalar-derivatives} gives, for sufficiently small
positive $s$,
\[
 \alpha_T(s)=
 \frac{\mathbb E V_{zzz}(s,Z_s)^2}
      {2\mathbb E V_{zz}(s,Z_s)^3}.
\]
Since $\alpha_T(s)\to0$, the diffusion has
$\mathbb E Z_s^2=Ts+o(s)$ and moments of order $2j$ of order $s^j$.
Oddness, the uniform spatial derivative bounds, and bounded time
derivatives imply
\[
 \mathbb E V_{zzz}(s,Z_s)^2
       =T V_{zzzz}(0,0)^2s+o(s),\qquad
 \mathbb E V_{zz}(s,Z_s)^3\longrightarrow V_{zz}(0,0)^3.
\]
The strict bounds in Lemma~\ref{prof:scalar-derivatives} yield
\[
 \kappa_T=\frac{T V_{zzzz}(0,0)^2}{2V_{zz}(0,0)^3}>0.
\]
Thus $q_T(u)\sim u/\kappa_T$ at zero. Full interval support makes
the distribution strictly increasing on its support and hence
makes the quantile continuous on $(0,1)$, even if the measure
has atoms. Away from zero, monotonicity, positivity, and
$q_T\le1$ extend the two-sided comparison with $u$ to the entire
open mass interval.
\end{proof}

\begin{proposition}[Rigidity at macroscopic scales]
\label{prof:macro-profile}
Consider minimizers in \eqref{eq:6} satisfying the endpoint and
constraint hypotheses \eqref{eq:3}--\eqref{eq:5}. Suppose
$x\ge c>0$, $e(x)/x^2\ge c>0$, and $d^2\le T-c$ along the
sequence. Then
\[
 r_n\longrightarrow q_T\quad\text{in }L^1(0,1)
 \quad\text{and locally uniformly on }(0,1).
\]
An omitted mass $u_0=o(a)$ and its specified admissible changed
paths are allowed, provided the changed paths of the minimizing
sequence and the smooth competitors used in the proof are admissible.
In particular all macroscopic quantile limits
are bounded above and below by fixed positive multiples of mass.
\end{proposition}

\begin{proof}
Pass to a subsequence on which $\lambda_n,d_n$ converge. The
assumed positive strength bounds force both limiting parameters
to be positive. Consequently $e_n$ converges uniformly on $[0,1]$
to a continuous strictly positive function $e$, with $T-e\ge c$.
The endpoint hypothesis \eqref{eq:4} implies
$b_0\to T$ and $h_0\to0$ in $L^1$ on the working masses. Indeed
$a=O(1)$, so each fixed positive mass satisfies $a\lesssim u$
with a fixed constant; boundedness then gives dominated convergence.
The omitted mass tends to zero and cannot affect these limits.

The constraints approximate every limiting quantile: their floor
is at most $n^{-3\delta}$ and their cap is at least $n^\delta$.
First approximate a desired distribution by a fixed smooth density
bounded above and below; its slight affine adjustment to the
endpoints $(r_{\min},u_0)$ satisfies the constraints for all
sufficiently large $n$. Here $r_{\min}\to0$ by \eqref{eq:5}.
The same construction supplies competitors on the working interval.
All changed paths used by the optimization and these competitors
are admissible by hypothesis.

Extract an $L^1$ limit $r$ of the minimizing quantiles. The uniform
pressure continuity bound in Lemma~\ref{prof:macro-variational}
and that lemma's variational limit, applied to the actual changed
paths, identify the limiting minimization as the joint minimum in
quantiles $r,z$ of
\[
 f_{\rm spin}\bigl(e(r)r+(T-e(r))z\bigr)
   +\frac14\int_0^1\bigl[e(r)r^2+(T-e(r))z^2\bigr]\,du.
\]
This reasoning applies both to the selected minimizing sequence
and to every approximating competitor, so it identifies a minimum
and its minimizers, not just an upper bound. It never evaluates a
possibly inadmissible bare endpoint.

Set
\[
 q=\frac{e(r)r+(T-e(r))z}{T}.
\]
This is a $[0,1]$-valued quantile: $re(r)$ is nondecreasing,
$T-e(r)$ is nonnegative and nondecreasing, and $z$ is
nonnegative and nondecreasing. Completing the square gives
\[
 e(r)r^2+(T-e(r))z^2
   =Tq^2+\frac{e(r)(T-e(r))}{T}(r-z)^2.
\]
The joint objective is therefore at least the ordinary scalar
minimum. Equality is attained by $r=z=q_T$. The coefficient of
$(r-z)^2$ is bounded below positively, so every minimizing pair
has $r=z=q$ almost everywhere. Lemma~\ref{prof:scalar-minimizer}
forces this quantile to be $q_T$.

Every subsequential limit is thus $q_T$, proving $L^1$ convergence
of the full sequence. Its continuity at each interior mass and
monotonicity of the approximating quantiles give local uniform
convergence. The artificial endpoint condition $r_n(1)=1$ is
irrelevant for this conclusion on the open mass interval.
\end{proof}
